
\section{The Newton Method\label{sec:Newton-Method}}

\begin{figure}
\centering{}\includegraphics[width=0.5\columnwidth]{Figures/Section_5\lyxdot 3}\caption{1D Newton Step}
\end{figure}

We begin with the Newton-Raphson method for finding the zeros of a
function $g$, i.e. finding $x$ such that $g(x)=0$. In optimization,
the goal is to find a root of the gradient, $\nabla f(x)=0$. The
Newton-Raphson iteration approximates a function by its linearization/Jacobian,
then guesses where the linear approximation intersects the zero line
as a likely zero, repeating this process until a sufficient approximation
has been found. In one dimension, we approximate the function $g$
by its linearization 
\[
g(x)\sim g(x^{(k)})+g'(x^{(k)})(x-x^{(k)}).
\]
Finding the zeros of the right-hand side, we have the Newton step
\[
x^{(k+1)}=x^{(k)}-\frac{g(x^{(k)})}{g'(x^{(k)})}.
\]
In high dimension, we can approximate the function by its Jacobian
$g(x)\sim g(x^{(k)})+Dg(x^{(k)})(x-x^{(k)})$ and this gives the step
\[
x^{(k+1)}=x^{(k)}-\left(Dg(x^{(k)})\right)^{-1}g(x^{(k)}).
\]

When $g(x)=\nabla f(x)$, then the Newton step becomes
\[
x^{(k+1)}=x^{(k)}-\left(\nabla^{2}f(x^{(k)})\right)^{-1}\nabla f(x^{(k)}).
\]

Alternatively, we can derive the Newton step as 
\[
x^{(k+1)}\leftarrow\arg\min_{x}f(x^{(k)})+\left\langle \nabla f(x^{(k)}),x-x^{(k)}\right\rangle +\frac{1}{2}(x-x^{(k)})^{\mathsf{T}}(\nabla^{2}f(x^{(k)}))(x-x^{(k)}).
\]
In words, the Newton step finds the minimum of the second order approximation
of the function. We note that the Newton iteration is a natural idea,
since
\begin{enumerate}
\item It is affine invariant (changing coordinate systems does not affect
convergence).
\item It is the fastest descent direction taking the Hessian into account.
\end{enumerate}
%
To see why affine-invariance is important for optimization, we consider
the following function
\[
f(x_{1},x_{2})=\frac{100}{2}x^{2}_{1}+\frac{1}{2}x^{2}_{2}.
\]
Gradient descent for this function is
\[
(x_{1},x_{2})\leftarrow(x_{1},x_{2})-h\nabla f(x_{1},x_{2})=((1-100h)x_{1},(1-h)x_{2}).
\]
We need $h<\frac{1}{100}$ in order for the first coordinate to converge,
but this will make the second coordinate converge too slowly. In general,
we may want to take different step sizes for different directions
and Newton method gives the best step if the function is quadratic.

For many classes of functions, gradient methods converge to the solution
linearly (namely, it takes $c\cdot\log\frac{1}{\epsilon}$ iterations
for some $c$ depending on the problem) while the Newton method converges
to the solution quadratically (namely, it takes $c'\cdot\log\log\frac{1}{\epsilon}$
for some $c'$) if the starting point is sufficiently close to a root.
However, each step of Newton method involves solving a linear system,
which can be much more expensive. Furthermore, Newton method may not
converge if the starting point is far away.

\begin{algorithm2e}[H]

\caption{$\mathtt{NewtonMethod}$}

\SetAlgoLined

\textbf{Input:} Function $g:\Rn\to\Rn$, initial point $x^{(0)}\in\Rn$,
and iteration count $T$

\For{$k=0,1,\cdots,T-1$}{

$x^{(k+1)}=x^{(k)}-\left(Dg(x^{(k)})\right)^{-1}g(x^{(k)}).$

}

\Return $x^{(T)}$.

\end{algorithm2e}
\begin{thm}[Quadratic convergence]
Assume that $g:\Rn\rightarrow\Rn$ is twice continuously differentiable.
Let $x^{(k)}$ be the sequence given by the Newton method. Suppose
that $x^{(k)}$ converges to some $x^{*}$ such that $g(x^{*})=0$
and $Dg(x^{*})$ is invertible. Then, for $k$ large enough, we have
\[
\|x^{(k+1)}-x^{*}\|\leq\|(Dg(x^{*}))^{-1}D^{2}g(x^{*})\|_{\op}\cdot\|x^{(k)}-x^{*}\|^{2}.
\]
\end{thm}

\begin{proof}
Let $e^{(k)}=x^{*}-x^{(k)}$. Then, we have that
\[
g(x^{*})=g(x^{(k)})+Dg(x^{(k)})[e^{(k)}]+\int^{1}_{0}(1-s)D^{2}g(x^{(k)}+se^{(k)})[e^{(k)},e^{(k)}]ds.
\]
Hence, we have
\[
0=Dg(x^{(k)})^{-1}g(x^{(k)})+x^{*}-x^{(k)}+\int^{1}_{0}(1-s)Dg(x^{(k)})^{-1}D^{2}g(x^{(k)}+se^{(k)})[e^{(k)},e^{(k)}]ds.
\]
Since $x^{(k+1)}=x^{(k)}-Dg(x^{(k)})^{-1}g(x^{(k)})$, we have
\[
x^{(k+1)}-x^{*}=\int^{1}_{0}(1-s)Dg(x^{(k)})^{-1}D^{2}g(x^{(k)}+se^{(k)})[e^{(k)},e^{(k)}]ds.
\]
So, we have
\begin{align*}
\|x^{(k+1)}-x^{*}\| & \leq\int^{1}_{0}(1-s)\|Dg(x^{(k)})^{-1}D^{2}g(x^{(k)}+se^{(k)})[e^{(k)},e^{(k)}]\|ds\\
 & \leq\int^{1}_{0}(1-s)\|Dg(x^{(k)})^{-1}D^{2}g(x^{(k)}+se^{(k)})\|_{\op}ds\cdot\|x^{(k)}-x^{*}\|^{2}.
\end{align*}
Since $x^{(k)}\rightarrow x^{*}$, we have that 
\[
Dg(x^{(k)})^{-1}D^{2}g(x^{(k)}+se^{(k)})\rightarrow Dg(x^{*})^{-1}D^{2}g(x^{*})
\]
and for large enough $k$, we can assume 
\[
\|Dg(x^{(k)})^{-1}D^{2}g(x^{(k)}+se^{(k)})\|_{\op}\le2\|Dg(x^{*})^{-1}D^{2}g(x^{*})\|_{\op}.
\]
The result follows.
\end{proof}
The argument above uses only that $g\in\mathcal{C}^{2}$ and does
not require convexity. Without further global assumptions on $g$,
the Newton method does not always converge to a root. The argument
above only shows that if the algorithm converges, then it converges
quadratically eventually. We call this ``local'' convergence since
it only gives a bound when $x^{(k)}$ is close enough to $x^{*}$.
In comparison, all earlier analyses in this book are about global
convergence, bounding the total number of iterations. In practice,
both analyses are important; global convergence makes sure the algorithm
is robust and local quadratic convergence makes sure the algorithm
converges to machine accuracy quickly enough. The local quadratic
convergence is particularly important for simple problems such as
computing $\sqrt{x}$.

\subsection*{Finding the largest root of a real-rooted polynomial}

In general, the Newton method does not converge globally. Here, we
give an application that does converge ``globally''.
\begin{thm}
Let $g(x)=\sum^{n}_{i=0}a_{i}x^{i}$ be a polynomial with only real
roots, and let $\lambda_{1}$ be its largest root. Suppose that $\lambda_{1}\leq x^{(0)}$.
Then, after $O(n\log(\frac{1}{\epsilon}))$ iterations, the Newton
method finds $x^{(k)}$ such that
\[
\lambda_{1}\leq x^{(k)}\leq\lambda_{1}+\epsilon\cdot(x^{(0)}-\lambda_{1}).
\]
\end{thm}

\begin{proof}
Since the roots of $g$ are real, we can write
\[
g(x)=b\cdot\prod^{n}_{i=1}(x-\lambda_{i})
\]
for some $b$ with $\lambda_{n}\leq\lambda_{n-1}\leq\cdots\leq\lambda_{1}$.
Then, we have that $g'(x)=b\cdot\sum_{i}\prod_{j\neq i}(x-\lambda_{j})$
and hence
\[
\frac{g(x)}{g'(x)}=\frac{1}{\sum_{i}\frac{1}{x-\lambda_{i}}}.
\]
For any $x>\lambda_{1}$, we have 
\begin{align*}
\frac{g(x)}{g'(x)} & \leq\frac{1}{\frac{1}{x-\lambda_{1}}}=x-\lambda_{1},\\
\frac{g(x)}{g'(x)}\geq & \frac{1}{\sum_{i}\frac{1}{x-\lambda_{1}}}=\frac{x-\lambda_{1}}{n}.\tag{{1}}
\end{align*}
Hence, by induction, we have that $x^{(k)}\geq\lambda_{1}$ during
every step $k$ and that
\[
x^{(k+1)}-\lambda_{1}\leq\left(1-\frac{1}{n}\right)(x^{(k)}-\lambda_{1})
\]
where we use $(1)$ to show the inductive step above.
\end{proof}
\begin{xca}
Consider the following iteration: 
\[
x^{(t+1)}=x^{(t)}-\frac{1}{n^{1/k}}\frac{g^{(k-1)}(x^{(t)})}{g^{(k)}(x^{(t)})}
\]
where $g^{(k)}(x)=\sum_{i}\frac{1}{(x-\lambda_{i})^{k}}$. For $k=1$,
this is exactly the Newton iteration as $g^{(0)}(x)=n$. Show that
when applied to a degree $n$ real-rooted polynomial, starting with
$x^{(0)}>\lambda_{1}$, in each iteration the distance to the largest
root decreases by a factor of $1-\frac{1}{n^{1/k}}$.
\end{xca}

Such a dependence of $\log\frac{1}{\epsilon}$ is called linear convergence.
The convergence of the Newton method can be quadratic, when close
enough to a root.
\begin{thm}
Assume that $f(x^{*})=0$, $|f'(x^{*})|\geq\alpha$, and $f'$ is
$L$-Lipschitz. Then, if $|x^{(0)}-x^{*}|\le\frac{\alpha}{2L}$,
\[
|x^{(k)}-x^{*}|\leq\frac{\alpha}{L}\left(\frac{L}{\alpha}|x^{(0)}-x^{*}|\right)^{2^{k}}.
\]
\end{thm}

\begin{proof}
The iteration is 
\[
x^{(k+1)}-x^{*}=x^{(k)}-x^{*}-\frac{f(x^{(k)})}{f'(x^{(k)})}.
\]
\begin{align*}
f(x^{*}) & =f(x^{(k)})+\int^{x^{*}}_{x^{(k)}}f'(z)dz=f(x^{(k)})+f'(x^{(k)})(x^{*}-x^{(k)})+\int^{x^{*}}_{x^{(k)}}(f'(z)-f'(x^{(k)}))dz.
\end{align*}
Therefore, 
\begin{align*}
\left|x^{(k+1)}-x^{*}\right| & =\frac{1}{|f'(x^{(k)})|}\left|\int^{x^{*}}_{x^{(k)}}(f'(z)-f'(x^{(k)}))dz\right|\\
 & \leq\frac{L|x^{*}-x^{(k)}|^{2}}{2|f'(x^{(k)})|}.
\end{align*}
Since $|f'(x^{(k)})|\geq|f'(x^{*})|-L|x^{(k)}-x^{*}|$
\[
\left|x^{(k+1)}-x^{*}\right|\le\frac{L}{2(\alpha-L|x^{(k)}-x^{*}|)}|x^{(k)}-x^{*}|^{2}.
\]
So, if $|x^{(k)}-x^{*}|\leq\frac{\alpha}{2L}$,
\[
|x^{(k+1)}-x^{*}|\leq\frac{L}{\alpha}|x^{(k)}-x^{*}|^{2}\le\frac{1}{2}|x^{(k)}-x^{*}|
\]
and 
\[
\frac{L}{\alpha}|x^{(k+1)}-x^{*}|\leq\left(\frac{L}{\alpha}|x^{(k)}-x^{*}|\right)^{2}.
\]
Writing $\epsilon_{0}=|x^{(0)}-x^{*}|$, after $k$ steps,
\[
\frac{L}{\alpha}|x^{(k)}-x^{*}|\leq\left(\frac{L}{\alpha}\epsilon_{0}\right)^{2^{k}}
\]
This implies $|x^{(k)}-x^{*}|<\epsilon$ after $O(\log\log(\frac{\alpha}{L\epsilon}))$
steps for fixed $\epsilon_{0}$. Note that $f$ has not been assumed
to be convex or polynomial.
\end{proof}
Moving back to optimization, we can view the goal as finding a root
of $\nabla f(x)=0$. Newton's iteration is the update 
\[
x^{(k+1)}=x^{(k)}-(\nabla^{2}f(x^{(k)}))^{-1}\nabla f(x^{(k)}).
\]
By the above proof, Newton's iteration has quadratic convergence from
points close enough to the optimum.

\subsection*{Quasi-Newton Method}

\begin{figure}
\centering{}\includegraphics[width=0.5\columnwidth]{Figures/Section_5\lyxdot 3(1)}\caption{Secant Method}
\end{figure}

When the Jacobian of $g$ is not available, we can approximate it
using the function itself. In one dimension, we can approximate the
Newton method and get the following secant method 
\[
x^{(k+1)}=x^{(k)}-g(x^{(k)})\frac{x^{(k)}-x^{(k-1)}}{g(x^{(k)})-g(x^{(k-1)})}
\]
where we approximate $g'(x^{(k)})$ by $\frac{g(x^{(k)})-g(x^{(k-1)})}{x^{(k)}-x^{(k-1)}}$.
For nice enough functions, the convergence rate satisfies $\varepsilon_{k+1}\leq C\cdot\varepsilon^{\frac{1+\sqrt{5}}{2}}_{k}$,
which is superlinear but not quadratic.

For $g$ that is high-dimensional (of dimension at least $2$), we
need to approximate its Jacobian. Let $J^{(k)}$ be the approximate
Jacobian we maintain in the $k^{\mathrm{th}}$ iteration. Similar
to the secant method, we want to enforce
\[
J^{(k+1)}\cdot(x^{(k)}-x^{(k-1)})=g(x^{(k)})-g(x^{(k-1)}).
\]
In dimension higher than one, this does not uniquely define the $J^{(k+1)}$.
One natural choice is to find $J^{(k+1)}$ that is closest to $J^{(k)}$
while satisfying the equation above. Solving the problem
\begin{equation}
\min_{J^{(k+1)}\cdot(x^{(k)}-x^{(k-1)})=g(x^{(k)})-g(x^{(k-1)})}\|J^{(k+1)}-J^{(k)}\|_{F},\label{eq:broyden_method_opt}
\end{equation}
we have the update rule
\begin{equation}
J^{(k+1)}=J^{(k)}+\frac{y^{(k)}-J^{(k)}s^{(k)}}{\|s^{(k)}\|^{2}}s^{(k)\top}\label{eq:broyden_method}
\end{equation}
where $s^{(k)}=x^{(k)}-x^{(k-1)}$ and $y^{(k)}=g(x^{(k)})-g(x^{(k-1)})$.
\begin{xca}
Prove that the equation (\ref{eq:broyden_method}) is indeed the minimizer
of (\ref{eq:broyden_method_opt}).
\end{xca}

When $g$ is given by the gradient of a convex function $f$, we know
that the Jacobian of $g$ satisfies $Dg=\nabla^{2}f(x)\succeq0$.
Therefore, we should impose some conditions such that $J^{(k)}\succeq0$
for all $k$.

\subsection*{BFGS algorithm}

Broyden--Fletcher--Goldfarb--Shanno (BFGS) algorithm is one of
the most popular quasi-Newton methods. The algorithm maintains an
approximate Hessian $J^{(k)}$ such that
\begin{itemize}
\item $J^{(k+1)}(x^{(k)}-x^{(k-1)})=\nabla f(x^{(k)})-\nabla f(x^{(k-1)})$
\item $J^{(k+1)}$ is close to $J^{(k)}$.
\item $J^{(k)}\succ0$.
\end{itemize}
To achieve all of these conditions, the natural optimization is 
\[
J^{(k+1)}\defeq\arg\min_{Js=y,J=J^{\top}}\norm{W^{\frac{1}{2}}\left(J^{-1}-\left(J^{(k)}\right)^{-1}\right)W^{\frac{1}{2}}}^{2}_{F}
\]
where $s^{(k)}=x^{(k)}-x^{(k-1)}$, $y^{(k)}=\nabla f(x^{(k)})-\nabla f(x^{(k-1)})$
and $W=\int^{1}_{0}\nabla^{2}f(x^{(k-1)}+s(x^{(k)}-x^{(k-1)}))ds$
(or any $W$ such that $Ws=y$). In some sense, the $W^{-\frac{1}{2}}$
is just a correct change of variables so that the algorithm is affine
invariant. Solving the equation above \cite{fletcher1987practical},
one obtains the update
\begin{equation}
\left(J^{(k+1)}\right)^{-1}=(I-\rho_{k}\cdot s^{(k)}y^{(k)\top})\left(J^{(k)}\right)^{-1}(I-\rho_{k}\cdot y^{(k)}s^{(k)\top})+\rho_{k}\cdot s^{(k)}s^{(k)\top}\label{eq:recursive}
\end{equation}
where $\rho_{k}=\frac{1}{y^{(k)\top}s^{(k)}}$. Alternatively, one
can also show that \cite{fletcher1991new}
\[
J^{(k+1)}=\arg\min_{Js=y,J=J^{\top}}D_{KL}(\mathcal{N}(0,J)\|\mathcal{N}(0,J^{(k)})).
\]

To implement the BFGS algorithm, it suffices to compute $\left(J^{(k)}\right)^{-1}\nabla f(x^{(k)})$.
Therefore, we can directly use the recursive formula (\ref{eq:recursive})
to compute $\left(J^{(k)}\right)^{-1}\nabla f(x^{(k)})$ instead of
maintaining $J^{(k)}$ or $\left(J^{(k)}\right)^{-1}$ explicitly.

In practice, the recursive formula for $\left(J^{(k)}\right)^{-1}$
becomes too expensive, and hence one can truncate the recursion after
a constant number of steps, which gives the limited-memory BFGS algorithm.

\subsection*{}
